
%\needspace{4\baselineskip}
\color{uniblue}{\subsection{Логика отключения трансформатора (ЛО Т)}}
\color{black}

\begin{enumerate}[label=\arabic{section}.\arabic{subsection}.\arabic*, labelsep=4pt, leftmargin=0pt, itemindent=57pt]

%%%%%%%%%%%%%%%%%%%%%%%%%%%%%%%%%%%%%%%%%%%%%%%%%%%%%%%%%%%%%%%%%%%%% ОБЩИЕ СВЕДЕНИЯ %%%%%%%%%%%%%%%%%%%%%%%%%%%%%%%%%%%%%%%%%%%%%%%%%%%%%%%%%%%%%%
\item Общие сведения
\begin{enumerate}[label=\arabic{section}.\arabic{subsection}.\arabic{enumi}.\arabic*, labelsep=4pt, leftmargin=0em, itemindent=65pt, parsep=0pt]

\item
Логика отключения трансформатора принимает сигналы срабатывания защитных функций трансформатора в составе устройства и формирует сигнал отключения трансформатора, сигнал запрета АПВ выключателя стороны НН трансформатора и сигнал запрета АВР секций НН. В устройстве логика отключения трансформатора представлена функциональным блоком «ЛО Т».

\item
В состав функционального блока «ЛО Т» входит непосредственно логика отключения (ЛО), логика запрета АПВ (ЗАПВ) и логика запрета АВР (ЗАВР).

%\item
%Функциональный блок <<ЛО Т>> принимает сигналы срабатывания защитных функций трансформатора в составе устройства и формирует сигнал отключения трансформатора, сигнал запрета АПВ выключателя стороны НН трансформатора и сигнал запрета АВР секций НН.

\item
Логика отключения трансформатора может быть выведена из работы оперативно путем активации сигнала «ОВ~ЛО~Т».

\item
Алгоритм работы ЛО Т выполнен в соответствии с рисунком \ref{loobsh:img1}. В таблице \ref{loobsh:tbl1} приведены параметры, необходимые для настройки ЛО Т.

% ЗДЕСЬ ВЕТВЛЕНИЕ ДВЕ ТАБЛИЦЫ ДЛЯ Т и Т2
\ifboolexpr{bool {isT2}}
{
\small
\begin{longtable}{|>{\centering\arraybackslash}m{5.3cm}|>{\centering\arraybackslash}m{3.3cm}|>{\centering\arraybackslash}m{4.2cm}|>{\centering\arraybackslash}m{1.8cm}|>{\centering\arraybackslash}m{1cm}|}
\caption{Параметры для настройки функции <<ЛО>>\hfill\vspace{-0.5\baselineskip}}\label{loobsh:tbl1}\\ 
\hline
\rowcolor{gray!30}
Параметр (Параметр на ИЧМ) & Условное обозначение на схеме & Значение/ Диапазон & Единица измерения & Шаг \\ 
\hline
\endfirsthead
\caption*{\hspace{3pt}\emph{Продолжение таблицы \ref{loobsh:tbl1}\hfill\vspace{-0.5\baselineskip}}} \\ % сделано по ГОСТ 2.105 п.6.8.7
\hline
\rowcolor{gray!30}
Параметр (Параметр на ИЧМ) & Условное обозначение на схеме & Значение/ Диапазон & Единица измерения & Шаг \\ 
%\hline
\endhead
%\hline
\endfoot
\hline
\endlastfoot
%==+t1*PTRC1|TRESOFFLVLGC
\multicolumn{5}{|c|}{ ЛО } \\ \hline 
\centering Ввод функции в работу (Ввод\_функции)  & \centering SGF1 & \centering Не предусмотрено \\ Предусмотрено & \centering -- & \centering \arraybackslash -- \\
\hline
\multicolumn{5}{|c|}{ ЗАПВ } \\ \hline 
\centering Ввод функции в работу (Ввод\_функции)  & \centering SGF1 & \centering Не предусмотрено \\ Предусмотрено & \centering -- & \centering \arraybackslash -- \\
\hline
\centering Режим запрета АПВ при срабатывании МТЗ 2 ступени (ЗАПВ\_МТЗ\_2ст)  & \centering SGF2 & \centering Не предусмотрено \\ Предусмотрено & \centering -- & \centering \arraybackslash -- \\
\hline
\centering Режим запрета АПВ при срабатывании МТЗ 3 ступени (ЗАПВ\_МТЗ\_3ст)  & \centering SGF3 & \centering Не предусмотрено \\ Предусмотрено & \centering -- & \centering \arraybackslash -- \\
\hline
\multicolumn{5}{|c|}{ ЗАВР } \\ \hline 
\centering Ввод функции в работу (Ввод\_функции)  & \centering SGF1 & \centering Не предусмотрено \\ Предусмотрено & \centering -- & \centering \arraybackslash -- \\
\hline
\centering Режим запрета АВР при срабатывании МТЗ 2 ступени (ЗАВР\_МТЗ\_2ст)  & \centering SGF2 & \centering Не предусмотрено \\ Предусмотрено & \centering -- & \centering \arraybackslash -- \\
\hline
\centering Режим запрета АВР при срабатывании МТЗ 3 ступени (ЗАВР\_МТЗ\_3ст)  & \centering SGF3 & \centering Не предусмотрено \\ Предусмотрено & \centering -- & \centering \arraybackslash -- \\
%===t1
\end{longtable}
\normalsize
}{}

\ifboolexpr{bool {isT}}
{
\small
\begin{longtable}{|>{\centering\arraybackslash}m{5.3cm}|>{\centering\arraybackslash}m{3.3cm}|>{\centering\arraybackslash}m{4.2cm}|>{\centering\arraybackslash}m{1.8cm}|>{\centering\arraybackslash}m{1cm}|}
\caption{Параметры для настройки функции <<ЛО>>\hfill\vspace{-0.5\baselineskip}}\label{loobsh:tbl1}\\ 
\hline
\rowcolor{gray!30}
Параметр (Параметр на ИЧМ) & Условное обозначение на схеме & Значение/ Диапазон & Единица измерения & Шаг \\ 
\hline
\endfirsthead
\caption*{\hspace{3pt}\emph{Продолжение таблицы \ref{loobsh:tbl1}\hfill\vspace{-0.5\baselineskip}}} \\ % сделано по ГОСТ 2.105 п.6.8.7
\hline
\rowcolor{gray!30}
Параметр (Параметр на ИЧМ) & Условное обозначение на схеме & Значение/ Диапазон & Единица измерения & Шаг \\ 
%\hline
\endhead
%\hline
\endfoot
\hline
\endlastfoot
%==+t1*PTRC1|TOFFLVLGC
\multicolumn{5}{|c|}{ ЛО } \\ \hline 
\centering Ввод функции в работу (Ввод\_функции)  & \centering SGF1 & \centering Не предусмотрено \\ Предусмотрено & \centering -- & \centering \arraybackslash -- \\
\hline
\multicolumn{5}{|c|}{ ЗАПВ } \\ \hline 
\centering Ввод функции в работу (Ввод\_функции)  & \centering SGF1 & \centering Не предусмотрено \\ Предусмотрено & \centering -- & \centering \arraybackslash -- \\
\hline
\centering Режим запрета АПВ при срабатывании МТЗ 2 ступени (ЗАПВ\_МТЗ\_2ст)  & \centering SGF2 & \centering Не предусмотрено \\ Предусмотрено & \centering -- & \centering \arraybackslash -- \\
\hline
\centering Режим запрета АПВ при срабатывании МТЗ 3 ступени (ЗАПВ\_МТЗ\_3ст)  & \centering SGF3 & \centering Не предусмотрено \\ Предусмотрено & \centering -- & \centering \arraybackslash -- \\
\hline
\multicolumn{5}{|c|}{ ЗАВР } \\ \hline 
\centering Ввод функции в работу (Ввод\_функции)  & \centering SGF1 & \centering Не предусмотрено \\ Предусмотрено & \centering -- & \centering \arraybackslash -- \\
\hline
\centering Режим запрета АВР при срабатывании МТЗ 2 ступени (ЗАВР\_МТЗ\_2ст)  & \centering SGF2 & \centering Не предусмотрено \\ Предусмотрено & \centering -- & \centering \arraybackslash -- \\
\hline
\centering Режим запрета АВР при срабатывании МТЗ 3 ступени (ЗАВР\_МТЗ\_3ст)  & \centering SGF3 & \centering Не предусмотрено \\ Предусмотрено & \centering -- & \centering \arraybackslash -- \\
%===t1
\end{longtable}
\normalsize
}{}

\end{enumerate}
%%%%%%%%%%%%%%%%%%%%%%%%%%%%%%%%%%%%%%%%%%%%%%%%%%%%%%%%%%%%%%%%%%%%% КОНЕЦ ОБЩИЕ СВЕДЕНИЯ %%%%%%%%%%%%%%%%%%%%%%%%%%%%%%%%%%%%%%%%%%%%%%%%%%%%%%%%%%%%%%


%%%%%%%%%%%%%%%%%%%%%%%%%%%%%%%%%%%%%%%%%%%%%%%%%%%%%%%%%%%%%%%% ЛОГИКА ОТКЛЮЧЕНИЯ ОБЩАЯ Т %%%%%%%%%%%%%%%%%%%%%%%%%%%%%%%%%%%%%%%%%%%%%%%%%%
\item Логика отключения (ЛО)
\begin{enumerate}[label=\arabic{section}.\arabic{subsection}.\arabic{enumi}.\arabic*, labelsep=4pt, leftmargin=0em, itemindent=65pt, parsep=0pt]

\item
Логика отключения в составе функционального блока «ЛО Т» представлена функцией «ЛО».

\item
Ввод функции в работу на этапе параметрирования устройства осуществляется программным переключателем SGF1 <<Ввод\_функции>>. Информация о введенной функции отображается выходным сигналом <<ЛО Т / ЛО: Ввод>>.
\item
Функция может быть выведена из работы оперативно путем активации сигнала <<ОВ ЛО Т / ЛО>>, что характеризуется активным состоянием сигнала <<ЛО Т / ЛО: Оперативный вывод>> на выходе функции.
\item
При появлении сигналов срабатывания от защитных функций в составе устройства функция формирует сигнал <<ЛО Т / ЛО: Срабатывание>>.

\FloatBarrier % Все float'ы ДО этой строки будут размещены здесь.

\end{enumerate}
%%%%%%%%%%%%%%%%%%%%%%%%%%%%%%%%%%%%%%%%%%%%%%%%%%%%%%%%%%%%%%%% КОНЕЦ ЛОГИКА ОТКЛЮЧЕНИЯ ОБЩАЯ Т %%%%%%%%%%%%%%%%%%%%%%%%%%%%%%%%%%%%%%%%%%%%%%%%%%


%%%%%%%%%%%%%%%%%%%%%%%%%%%%%%%%%%%%%%%%%%%%%%%%%%%%%%%%%%%%%%%% ЛОГИКА ЗАПРЕТА АПВ НН %%%%%%%%%%%%%%%%%%%%%%%%%%%%%%%%%%%%%%%%%%%%%%%%%%
\item Логика запрета АПВ выключателей трансформатора (ЗАПВ)
\begin{enumerate}[label=\arabic{section}.\arabic{subsection}.\arabic{enumi}.\arabic*, labelsep=4pt, leftmargin=0em, itemindent=65pt, parsep=0pt]

\item
Логика запрета АПВ выключателей трансформатора в составе функционального блока «ЛО Т» представлена функцией «ЗАПВ».

\item
Ввод функции в работу на этапе параметрирования устройства  осуществляется программным переключателем SGF1 <<Ввод\_функции>>. Информация о введенной функции отображается выходным сигналом <<ЛО Т / ЗАПВ: Ввод>>.
\item
Функция может быть выведена из работы оперативно путем активации сигнала <<ОВ ЛО Т / ЗАПВ>>, что характеризуется активным состоянием сигнала <<ЛО Т / ЗАПВ: Оперативный вывод>> на выходе функции.
\item
Функция принимает сигналы срабатывания от защит в составе устройства (сигналы срабатывания от МТЗ 2 и 3 ступеней заводятся через программные переключатели SGF2 <<ЗАПВ\_МТЗ\_ст2>> и SGF3 <<ЗАПВ\_МТЗ\_ст3>> соответственно) и формирует сигнал <<ЛО Т / ЗАПВ: Запрет АПВ>>.

\end{enumerate}

\FloatBarrier % Все float'ы ДО этой строки будут размещены здесь.


%%%%%%%%%%%%%%%%%%%%%%%%%%%%%%%%%%%%%%%%%%%%%%%%%%%%%%%%%%%%%%%% ЛОГИКА ЗАПРЕТА АВР НН %%%%%%%%%%%%%%%%%%%%%%%%%%%%%%%%%%%%%%%%%%%%%%%%%%
\item Логика запрета АРВ секций НН (ЗАВР)
\begin{enumerate}[label=\arabic{section}.\arabic{subsection}.\arabic{enumi}.\arabic*, labelsep=4pt, leftmargin=0em, itemindent=65pt, parsep=0pt]

\item
Логика запрета АВР секций НН в составе функционального блока «ЛО Т» представлена функцией «ЗАВР».

\item
Ввод функции в работу на этапе параметрирования устройства  осуществляется программным переключателем SGF1 <<Ввод\_функции>>. Информация о введенной функции отображается выходным сигналом <<ЛО Т / ЗАВР: Ввод>>.
\item
Функция принимает сигналы срабатывания от МТЗ 2 и 3 ступеней (через программные переключатели SGF2 <<ЗАВР\_МТЗ\_ст2>> и SGF3 <<ЗАВР\_МТЗ\_ст3>> соответственно) и формирует сигнал <<ЛО Т / ЗАВР: Запрет АВР>>. 

\end{enumerate}

%%%%%%%%%%%%%%%%%%%%%%%%%%%%%%%%%%%%%%%%%%%%%%%%%%%%%%%%%%%%%%%%%%%%%%%
\vspace{1mm}
\begin{figure}[H]
\centering
% ЗДЕСЬ ВЕТВЛЕНИЕ, рисунок для исполнения Т2
\ifboolexpr{bool {isT2}}
{\includegraphics[width=1\textwidth,height=1\textheight,keepaspectratio]{\fbpath/35. ЛО Т 35 Т/_latex/img/_2.pdf}}{}

% ЗДЕСЬ ВЕТВЛЕНИЕ, рисунок для исполнения Т
\ifboolexpr{bool {isT}}
{\includegraphics[width=1\textwidth,height=1\textheight,keepaspectratio]{\fbpath/35. ЛО Т 35 Т/_latex/img/_1.pdf}}{}

\captionof{figure}{Функциональная схема ЛО Т}\label{loobsh:img1}
\end{figure}
%%%%%%%%%%%%%%%%%%%%%%%%%%%%%%%%%%%%%%%%%%%%%%%%%%%%%%%%%%%%%%%%%%%%%%%%

\FloatBarrier % Все float'ы ДО этой строки будут размещены здесь.

\end{enumerate}